Robots navigating in human-populated environments must not only avoid collisions but also respect the social structure of human crowds, particularly the implicit boundaries of social groups. Most existing navigation approaches model humans as independent individuals, leading to socially disruptive behavior even when collision-free motion is achieved. This paper presents TAGA (Tangent Action for Group Avoidance), a modular and lightweight framework that explicitly accounts for human group formations during robot navigation. TAGA detects groups online from observable motion cues using a DBSCAN-based module and applies a tangent-based action strategy to guide the robot smoothly around group boundaries while maintaining efficient progress toward the goal. A hierarchical safety controller coordinates group-level avoidance with individual collision prevention. We further propose the Group Collision Rate (GCR) as a continuous complementary metric that quantifies group boundary violations throughout the navigation episode, providing finer-grained assessment than terminal-state metrics alone. To support rigorous evaluation, we introduce a realistic crowd simulation benchmark incorporating five empirically motivated modeling phases: individual speed variation, dynamic group speed coupling, F-formations for static groups, leader--follower dynamics, and convex-hull group boundaries. Extensive experiments across reactive methods (ORCA, Social Force) and learning-based methods (DS-RNN, AG-RL, GARN) reveal a consistent pattern: TAGA significantly improves social compliance for reactive methods while the benefit diminishes for highly optimized learning-based policies that have already internalized implicit group-avoidance behavior through training. These findings provide practical guidance for when modular group-awareness adds the most value, and demonstrate TAGA as an effective plug-in solution for socially compliant robot navigation in realistically modeled crowd environments.

% Robots operating in human-populated environments must navigate safely while also respecting social norms, particularly those associated with human group behavior. Most existing navigation methods focus on avoiding individual humans and often overlook group structures, which can lead to socially disruptive behavior despite maintaining collision-free motion. This paper introduces TAGA, a modular and lightweight framework for socially aware robot navigation that explicitly accounts for human group formations. TAGA detects groups online using motion-based cues and applies a tangent-based action strategy that enables the robot to bypass group boundaries smoothly while maintaining efficient progress toward the goal. To ensure both social and physical safety, a hierarchical safety controller is designed to coordinate group-level avoidance with individual collision prevention. In addition, we propose the Group Collision Rate, a continuous metric that measures the extent to which a robot violates group boundaries during navigation. Extensive simulation experiments across dense, dynamic, and mixed crowd scenarios demonstrate that TAGA significantly reduces group intrusions while preserving or improving navigation success rates. The results show that TAGA can be seamlessly integrated with existing geometric and learning-based navigation systems, offering a practical solution for socially compliant robot navigation in real-world environments.
% \SR{Robot navigation in densely populated environments presents significant challenges, particularly regarding the interplay between individual and group dynamics. Current navigation models predominantly address interactions with individual pedestrians while failing to account for human groups that naturally form in real-world settings. Conversely, the limited models implementing group-aware navigation typically prioritize group dynamics at the expense of individual interactions, both of which are essential for socially appropriate navigation. This research extends an existing simulation framework to incorporate both individual pedestrians and human groups. We present Tangent Action for Group Avoidance (TAGA), a modular reactive mechanism that can be integrated with existing navigation frameworks to enhance their group-awareness capabilities. TAGA dynamically modifies robot trajectories using tangent action-based avoidance strategies while preserving the underlying model's capacity to navigate around individuals. Additionally, we introduce Group Collision Rate (GCR), a novel metric to quantitatively assess how effectively robots maintain group integrity during navigation. Through comprehensive simulation-based benchmarking, we demonstrate that integrating TAGA with state-of-the-art navigation models (ORCA, Social Force, DS-RNN, and AG-RL) reduces group intrusions by 45.7–78.6\% while maintaining comparable success rates and navigation efficiency. Future work will focus on real-world implementation and validation of this approach. 



% Navigating safely and efficiently in dense human crowds remains a critical challenge for autonomous robots. Traditional navigation models primarily focus on individual human interactions, often neglecting human groups that naturally form in real-world environments. While some models incorporate group-aware navigation, they primarily focus on group dynamics while overlooking interactions with individual pedestrians, which are equally crucial for socially compliant navigation. To address this, we extend an existing simulation framework to include both individual humans and human groups. Additionally, we introduce Group Collision Rate (GCR), a novel metric to quantify how well a robot respects group integrity. Building on this, we propose Tangent Action for Group Avoidance (TAGA), a modular and reactive group-aware navigation mechanism that dynamically adjusts robot trajectories using a tangent action-based avoidance strategy while seamlessly integrating with individual-based navigation models. .Finally, we benchmark our approach by integrating TAGA with state-of-the-art navigation models (ORCA, Social Force, DS-RNN, and AG-RL) and comparing their performance with and without TAGA enhancement. The results show that TAGA reduces group intrusions by 45.7–78.6\% while maintaining comparable success rates and navigation efficiency.
% % Finally, we benchmark TAGA against state-of-the-art models (ORCA, Social Force, DS-RNN, and AG-RL). The results show that TAGA reduces group intrusions by  45.7–78.6\% while maintaining comparable success rates and navigation efficiency. 
% Additional details and resources can be found at: \url{https://sites.google.com/rme.du.ac.bd/taga/home}.
% Navigating safely and efficiently in dense human crowds remains a critical challenge for autonomous robots. Traditional navigation models primarily focus on individual human interactions, often failing to account for human groups that naturally form in real-world environments. While some models incorporate group-aware navigation, they primarily focus on group dynamics, overlooking interactions with individual pedestrians, which are equally crucial for socially compliant navigation. In this work, we introduce Tangent Action for Group Avoidance (TAGA), a reactive group-aware navigation model that enables robots to dynamically adapt their trajectories based on detected human groups using a tangent action-based avoidance strategy, while seamlessly integrating with individual-based navigation models. By extending an existing simulation framework, we incorporate realistic group behaviors and evaluate TAGA against state-of-the-art navigation models. Additionally, we propose Group Collision Rate (GCR), a novel performance metric that quantifies a robot’s ability to respect group integrity. Benchmarking results demonstrate that TAGA significantly reduces group intrusions while maintaining competitive navigation efficiency. Our approach offers a practical solution for improving socially compliant robot navigation in dynamic, interactive human environments. Our project website with additional details and resources can be found at: \url{https://sites.google.com/rme.du.ac.bd/taga/home}.

% This electronic document is a ÒliveÓ template. The various components of your paper [title, text, heads, etc.] are already defined on the style sheet, as illustrated by the portions given in this document.
